\documentclass[12pt,a4paper]{article}

% Paquetes de idioma y codificación
\usepackage[utf8]{inputenc}
\usepackage[spanish]{babel}

% Paquetes de diseño y formato
\usepackage[margin=2.5cm]{geometry}
\usepackage{listings}
\usepackage{xcolor}
\usepackage{hyperref}
\usepackage{titlesec}
\usepackage{graphicx}

% Colores personalizados para el resaltado YAML
\definecolor{yamlkey}{RGB}{0, 102, 204}
\definecolor{yamlval}{RGB}{204, 0, 0}
\definecolor{yamlcomment}{RGB}{120, 120, 120}
\definecolor{bglight}{RGB}{245, 245, 245}

% Configuración del lenguaje YAML para el paquete listings
\lstdefinelanguage{yaml}{
  keywords={nodes, materials, elements, mesh, mesh_material, mesh_thickness, boundary_conditions, boundary_conditions_by_coord, boundary_conditions_by_group, point_loads, point_loads_by_coord, point_loads_by_group, solver, output, file_name, pre_process, export, results, frequency, nodal_results, element_results, text_export, format, nodes, elements, id, type, E, nu, sigma_y, H, hypothesis, ux, uy, tol, coord, val, num_steps, max_iter, adaptive},
  keywordstyle=\color{yamlkey}\bfseries,
  ndkeywords={true, false, null},
  ndkeywordstyle=\color{yamlval}\bfseries,
  identifierstyle=\color{black},
  sensitive=false,
  comment=[l]{\#},
  commentstyle=\color{yamlcomment}\ttfamily,
  stringstyle=\color{yamlval}\ttfamily,
  morestring=[b]',
  morestring=[b]"
}

\lstset{
  language=yaml,
  basicstyle=\ttfamily\small,
  backgroundcolor=\color{bglight},
  frame=single,
  rulecolor=\color{lightgray},
  breaklines=true,
  showstringspaces=false
}

\title{\Huge \textbf{Manual de Referencia: Formato YAML} \\ \vspace{0.5cm} \Large Motor Fenix FEM}
\author{Documentación Oficial del Software}
\date{\today}

\begin{document}

\maketitle

\tableofcontents
\vspace{1cm}

\begin{abstract}
El presente documento detalla la estructura y todas las características implementadas en el analizador de archivos YAML (\texttt{YamlParser}) del software Fenix FEM. 

\textbf{Regla de Oro del Software:} Debe quedar perfectamente claro que \textbf{toda la información} necesaria para definir y resolver un problema en particular (geometría, materiales, apoyos, cargas y opciones de exportación) debe estar contenida exclusivamente en su respectivo archivo de entrada en formato YAML. \textbf{Bajo ninguna circunstancia se debe modificar el código fuente del programa} para correr un problema específico. Esta arquitectura estrictamente \textit{Data-Driven} garantiza que el motor matemático se mantenga universal e intocable.
\end{abstract}

\section{Definición de la Malla}
La geometría del dominio puede declararse de dos maneras distintas: importándola desde un generador de mallas externo (Gmsh) o declarándola nodo por nodo manualmente.

\subsection{Importación desde Gmsh}
Permite cargar mallas de geometrías complejas (archivos \texttt{.msh} en formato ASCII 4.1). Se debe especificar el archivo, el ID del material a asignar y el espesor transversal de la placa.
\begin{lstlisting}
mesh: "placa.msh"
mesh_material: 1
mesh_thickness: 0.05
\end{lstlisting}

\subsection{Definición Manual de Nodos y Elementos}
Ideal para armaduras o placas sencillas. Los nodos requieren su ID y coordenadas \texttt{[x, y]}. Los elementos requieren \texttt{[Tipo, ID\_Material, Espesor, Nodo1, Nodo2, ...]}. Tipos soportados: \texttt{Quad4}, \texttt{Tri3} y \texttt{Truss2D}.
\begin{lstlisting}
nodes:
  1: [0.0, 0.0]
  2: [1.0, 0.0]

elements:
  1: [Quad4, 1, 0.1, 1, 2, 5, 4]
  2: [Tri3,  1, 0.1, 2, 3, 6]
\end{lstlisting}

\section{Modelos de Materiales}
Se listan en el bloque \texttt{materials}. Actualmente Fenix FEM soporta materiales linealmente elásticos (\texttt{Elastic2D}) y elastoplásticos de Von Mises con endurecimiento (\texttt{VonMises2D}).
\begin{lstlisting}
materials:
  - id: 1
    type: Elastic2D
    E: 200.0e9
    nu: 0.3
    hypothesis: plane_stress
  - id: 2
    type: VonMises2D
    E: 200.0e9
    nu: 0.3
    sigma_y: 250.0e6
    H: 2.0e9
\end{lstlisting}

\section{Condiciones de Frontera y Cargas}
Fenix FEM posee una arquitectura completamente simétrica para la asignación de restricciones y fuerzas. \textbf{Cualquier método de selección de nodos} (por ID, coordenadas, o grupos físicos) aplica exactamente igual para imponer apoyos o desplazamientos prescritos (\texttt{boundary\_conditions}) que para aplicar fuerzas puntuales externas (\texttt{point\_loads}).

\subsection{Asignación Nodo por Nodo}
Se especifica el ID del nodo y el Grado de Libertad (\texttt{ux}, \texttt{uy}) a modificar.
\begin{lstlisting}
boundary_conditions:
  1: {ux: 0.0, uy: 0.0}   # Empotramiento

point_loads:
  3: {uy: -5000.0}        # Fuerza de 5 kN hacia abajo
\end{lstlisting}

\subsection{Búsqueda por Coordenadas Extremas (Ideal para Gmsh)}
Cuando se importan mallas desde Gmsh, generalmente se desconoce qué ID de número le corresponde a cada nodo. Para solucionar esto, Fenix FEM permite aplicar apoyos o fuerzas directamente a los \textbf{bordes físicos} de la geometría.

El motor matemático calcula la "caja" que encierra al modelo y te permite llamar a sus cuatro caras exteriores:
\begin{itemize}
    \item \texttt{x\_min}: Borde extremo izquierdo.
    \item \texttt{x\_max}: Borde extremo derecho.
    \item \texttt{y\_min}: Borde extremo inferior.
    \item \texttt{y\_max}: Borde extremo superior.
\end{itemize}

El programa buscará automáticamente a \textbf{todos los nodos} que se encuentren tocando ese borde y les aplicará el valor especificado. 

\textbf{Nota Importante:} Los valores declarados dentro de las llaves (ej. \texttt{ux: 0.0}) \textbf{no son coordenadas geométricas}, sino que representan el \textbf{Desplazamiento} o la \textbf{Fuerza} que se le va a aplicar a los nodos de dicho borde. Dado que las coordenadas manejan múltiples decimales, se exige el uso del parámetro \texttt{tol} (Tolerancia geométrica, ej. \texttt{1.0e-6}) para crear un pequeño margen de búsqueda que atrape a los nodos correctamente compensando los errores de redondeo.

\begin{lstlisting}
boundary_conditions_by_coord:
  # Empotra absolutamente todos los nodos ubicados en el borde izquierdo
  x_min: {ux: 0.0, uy: 0.0, tol: 1.0e-6}

point_loads_by_coord:
  # Aplica una fuerza de 15 kN hacia la derecha a cada nodo del borde derecho
  x_max: {ux: 15000.0, tol: 1.0e-6}
\end{lstlisting}

\subsection{Búsqueda por Coordenada Exacta (Láser)}
Si necesitas imponer un apoyo o carga en un punto interior de la malla (que no se encuentra en las fronteras externas), puedes crear un identificador personalizado indicando el eje de corte (\texttt{coord: 'x'} o \texttt{coord: 'y'}) y el valor exacto (\texttt{val}).
\begin{lstlisting}
boundary_conditions_by_coord:
  # Empotra todos los nodos que se encuentren exactamente en la linea x = 2.5
  apoyo_central: {coord: 'x', val: 2.5, uy: 0.0, tol: 1.0e-5}

point_loads_by_coord:
  # Aplica una fuerza en los nodos ubicados a 5 metros de altura
  fuerza_media: {coord: 'y', val: 5.0, uy: -5000.0, tol: 1.0e-5}
\end{lstlisting}

\subsection{Asignación por Grupos Físicos de Gmsh (Recomendado)}
La manera más robusta y profesional de aplicar condiciones de frontera es a través de los \textit{Physical Groups} de Gmsh. Al asignar un nombre a una línea o superficie en el modelo CAD, Fenix FEM puede atrapar automáticamente todos los nodos pertenecientes a dicho grupo.

\begin{lstlisting}
boundary_conditions_by_group:
  "Apoyo_Izquierdo": {ux: 0.0, uy: 0.0}

point_loads_by_group:
  "Carga_Motor": {uy: -5000.0}
\end{lstlisting}

\section{Configuración del Motor Matemático}
El bloque \texttt{solver} dictamina el comportamiento algorítmico de la solución. Puede ser \texttt{LinearSolver} para problemas elásticos directos de 1 paso, o \texttt{NonlinearSolver} para problemas inelásticos.
\begin{lstlisting}
solver:
  type: NonlinearSolver
  num_steps: 10       # Cantidad de incrementos de carga
  max_iter: 15        # Limite iterativo Newton-Raphson
  tol: 1.0e-5         # Tolerancia de la norma del error
  adaptive: true      # Control de paso mediante biseccion
\end{lstlisting}

\section{Opciones de Salida (Output)}
El bloque \texttt{output} le otorga al usuario el control absoluto sobre qué archivos debe generar el motor al finalizar la lectura y el cálculo. Esto asegura que el código fuente de ejecución jamás deba modificarse para habilitar o deshabilitar exportaciones.
\begin{lstlisting}
output:
  file_name: "mi_analisis"
  pre_process:
    export: true                # Genera archivo inicial de apoyos y cargas
  results:
    frequency: "all_steps"      # Opciones: "last_step" o "all_steps" (Animacion)
    nodal_results:              # Variables exportadas en los nodos
      - "Displacements"
    element_results:            # Variables exportadas en los elementos
      - "Von_Mises"
      - "Internal_State"
\end{lstlisting}

\subsection{Extracción de Datos Tabulares (Exportación a Texto)}
Adicionalmente, se puede colocar un bloque \texttt{text\_export} separado del bloque de resultados para extraer las variables de un grupo específico de nodos y elementos en formato de texto tabulado (TXT). Esto genera una tabla con columnas alineadas, ideal para lectura directa o para graficar curvas.
\begin{lstlisting}
  text_export:
    export: true
    format: "txt"
    nodes: "all"             # Puede ser una lista ej. [2, 15] o la palabra "all"
    nodal_results:
      - "Displacements"
    elements: [41]           # Puede ser una lista o la palabra "all"
    element_results:
      - "Von_Mises"
\end{lstlisting}

\end{document}